\documentclass{article}

% NeurIPS 2026 style configuration
\usepackage[preprint]{neurips_2026}

\usepackage[utf8]{inputenc} % allow utf-8 input
\usepackage[T1]{fontenc}    % use 8-bit T1 fonts
\usepackage{hyperref}       % hyperlinks
\usepackage{url}            % simple URL typesetting
\usepackage{booktabs}       % professional-quality tables
\usepackage{amsfonts}       % blackboard math symbols
\usepackage{nicefrac}       % compact symbols for 1/2, etc.
\usepackage{microtype}      % microtypography
\usepackage{xcolor}         % colors
\usepackage{graphicx}
\usepackage{amsmath,amssymb}

\title{Representation Regularization in Joint-Embedding Predictive Architectures for Multi-Modal Brain MRI: Label Efficiency and Out-of-Distribution Generalization}

\author{%
  Han Riman \thanks{Corresponding author.} \\
  Master's Thesis Research Laboratory \\
  \texttt{hanriman@master.thesis} \\
}

\begin{document}

\maketitle

\begin{abstract}
Joint-Embedding Predictive Architectures (JEPA) have emerged as a powerful self-supervised representation learning paradigm operating in latent embedding space without pixel-level reconstruction. Standard I-JEPA prevents representation collapse using an Exponential Moving Average (EMA) teacher target encoder heuristic. Recent statistical regularization advances—specifically SIGReg (Sketched Isotropic Gaussian Regularization in LeJEPA) and VISReg (Variance-Invariance-Sketching Regularization)—enable \emph{heuristic-free} single-encoder JEPA training without EMA teachers. In this paper, we systematically evaluate I-JEPA, SigReg JEPA, VisReg JEPA, a Residual UNet baseline, and a state-of-the-art nnU-Net baseline with deep supervision on 2D multi-modal Glioma MRI (BraTS GLI) across three key clinical paradigms: (1) Full-Data Benchmark, (2) Low-Data Label Efficiency ($1\%$ to $100\%$ labels), and (3) Out-of-Distribution (OOD) Scanner Shift and Cross-Pathology Transfer on Meningioma MRI (BraTS-MEN-RT). Our empirical findings demonstrate that at 10\% labels ($633$ training slices), SigReg JEPA achieves higher downstream Test Dice ($0.5198$) than nnU-Net ($0.5163$) with a $2.4\times$ lower Hausdorff boundary error ($16.44\,\text{px}$ vs $40.26\,\text{px}$). Furthermore, on zero-shot cross-pathology evaluation under zero-padded missing modalities, SigReg JEPA maintains top transferability ($0.1026$ Dice) while standard UNet completely collapses ($0.0078$ Dice), proving that analytical variance-covariance regularization creates robust, modality-invariant representations without momentum teacher overhead.
\end{abstract}


\section{Introduction}

Automated segmentation of brain gliomas from multi-modal Magnetic Resonance Imaging (MRI)—comprising T1, T1-post contrast (T1c), T2, and Fluid Attenuated Inversion Recovery (FLAIR) sequences—is fundamental for clinical neuro-oncology \citep{brats2024}. Fully supervised deep convolutional networks, such as Residual UNet \citep{ronneberger2015u} and nnU-Net \citep{isensee2021nnunet}, achieve strong performance when extensive pixel-wise ground-truth annotations are available. However, expert manual delineation is labor-intensive, rare, and susceptible to inter-observer variability.

Self-Supervised Learning (SSL) pre-trains representation encoders on unannotated MRI slices prior to downstream task fine-tuning. Unlike pixel-reconstructing Masked Autoencoders (MAE), Joint-Embedding Predictive Architectures (JEPA) \citep{assran2023ijepa} learn by predicting target patch representations from context patch representations in latent space. Standard I-JEPA relies on an Exponential Moving Average (EMA) teacher update to prevent feature collapse. Recent formulations, including SIGReg \citep{balestriero2025lejepa} and VISReg \citep{wu2026visreg}, introduce analytical feature regularization penalties that eliminate EMA teachers, establishing a \emph{heuristic-free} single-encoder framework.

In this work, we present a systematic comparative investigation of self-supervised JEPAs and supervised baselines across three clinical regimes:
\begin{enumerate}
    \item \textbf{Full-Data Benchmark}: Evaluating latent feature rank, cosine similarity, and downstream segmentation overlap.
    \item \textbf{Low-Data Label Efficiency ($1\%$ to $100\%$ labels)}: Evaluating model transferability when annotated training scans are extremely scarce ($11$ to $6,330$ slices).
    \item \textbf{Out-of-Distribution (OOD) Scanner Shift \& Cross-Pathology Transfer}: Evaluating resilience to synthetic scanner perturbations (Rician Noise SNR degradation, $B_1$ Bias Field coil shifts) and zero-shot transfer to single-sequence Extra-axial Meningioma MRI (\texttt{BraTS-MEN-RT}).
\end{enumerate}


\section{Methodology}

\subsection{Multi-Modal Input Formulation \& Random Modality Dropout}
Each 2D axial MRI slice $X \in \mathbb{R}^{B \times 4 \times 240 \times 240}$ encapsulates four co-registered anatomical sequences: $X = [X_{\text{T1}}, X_{\text{T1c}}, X_{\text{T2}}, X_{\text{FLAIR}}]$. Non-zero brain tissue voxels are Z-score normalized independently per volume. To prevent modality co-dependence, we inject \textbf{Random Modality Dropout} ($p_{\text{drop}} = 0.25$) during training:
\begin{equation}
    \mathbf{M}_{b, c} \sim \text{Bernoulli}(1 - p_{\text{drop}}), \qquad \hat{X}_{b, c, :, :} = X_{b, c, :, :} \odot \mathbf{M}_{b, c}
\end{equation}

\subsection{Vision Transformer Backbone \& Segmentation Decoder}
The vision encoder $E_\theta$ processes 2D multi-modal slices using $16 \times 16$ patch embeddings, mapping $240 \times 240$ slices to a grid of $N = 225$ patch tokens ($D = 384$). For downstream segmentation, a 4-stage upsampling decoder $D_\phi$ reshapes patch tokens to $[B, 384, 15, 15]$ and progressively upsamples features through transpose-convolutional stages ($15\times15 \rightarrow 30\times30 \rightarrow 60\times60 \rightarrow 120\times120 \rightarrow 240\times240$) to output 2D tumor segmentation logits $\hat{Y} \in \mathbb{R}^{B \times 1 \times 240 \times 240}$.

\subsection{JEPA Variants \& Regularization Objectives}

\paragraph{I-JEPA (EMA Teacher Dual-Encoder)}
Predicts target patch representations using student context encoder $E_\theta$, predictor $P_\phi$, and momentum target teacher encoder $E_{\bar{\theta}}$ updated at step $t$ via momentum $m = 0.996$ \citep{assran2023ijepa}:
\begin{equation}
    \bar{\theta}_t \leftarrow m \bar{\theta}_{t-1} + (1 - m) \theta_t
\end{equation}

\paragraph{SigReg JEPA (Heuristic-Free Single-Encoder)}
Eliminates EMA target encoders using Sketched Isotropic Gaussian Regularization (SIGReg) \citep{balestriero2025lejepa}:
\begin{equation}
    \mathcal{L}_{\text{SigReg}} = \mathcal{L}_{\text{I-JEPA}} + \lambda_{\text{var}} \mathcal{L}_{\text{var}}(Z) + \lambda_{\text{cov}} \mathcal{L}_{\text{cov}}(Z)
\end{equation}
where $\mathcal{L}_{\text{var}}$ forces feature dimension variance above $\gamma = 1.0$, and $\mathcal{L}_{\text{cov}}$ penalizes off-diagonal feature covariance.

\paragraph{VisReg JEPA (Spatial Variance Regularization)}
Applies spatial patch feature variance contrast penalties across spatial patch dimensions \citep{wu2026visreg}:
\begin{equation}
    \mathcal{L}_{\text{VisReg}} = \mathcal{L}_{\text{I-JEPA}} + \lambda_{\text{vis}} \cdot \frac{1}{B} \sum_{b=1}^B \max\left(0, \, 1.0 - \sqrt{\text{Var}_{\text{patch}}(Z_b) + \epsilon}\right)
\end{equation}


\section{Experiments and Empirical Results}

\subsection{Full-Data Benchmark Results}
Table~\ref{tab:benchmark_results} summarizes full-data pre-training and downstream segmentation performance. SigReg JEPA achieves near-maximal Effective Feature Rank ($\mathbf{366.91}$ out of $384$), low cosine similarity ($0.0309$), and top self-supervised downstream Dice ($0.8530$).

\begin{table}[ht]
\centering
\caption{Full-data empirical benchmark comparison of throughput, latency, representation collapse probes, and downstream segmentation metrics on 2D multi-modal BraTS Glioma MRI.}
\label{tab:benchmark_results}
\vspace{0.5em}
\resizebox{\textwidth}{!}{%
\begin{tabular}{lccccccc}
\toprule
\textbf{Model Architecture} & \textbf{Test Dice} $\uparrow$ & \textbf{Test IoU} $\uparrow$ & \textbf{HD95 (px)} $\downarrow$ & \textbf{Effective Rank} $\uparrow$ & \textbf{Avg Cosine Sim} $\downarrow$ & \textbf{Inference Latency} & \textbf{Training Speed} \\
\midrule
\textbf{UNet Baseline} & 0.8691 & 0.7701 & 5.22 px & N/A (CNN) & N/A & \textbf{19.18 ms/slice} & \textbf{17.61 s/epoch} \\
\textbf{nnU-Net Baseline (SOTA)} & \textbf{0.9073} & \textbf{0.8319} & \textbf{3.34 px} & N/A (CNN) & N/A & 20.59 ms/slice & 33.65 s/epoch \\
\textbf{I-JEPA (Fine-tuned)} & 0.8437 & 0.7332 & 11.06 px & 124.98 & 0.9556 & 18.40 ms/slice & 20.32 s/epoch \\
\textbf{SigReg JEPA (Fine-tuned)} & \textbf{0.8530} & \textbf{0.7464} & \textbf{10.75 px} & \textbf{366.91} (Max) & \textbf{0.0309} & \textbf{18.38 ms/slice} & \textbf{20.16 s/epoch} \\
\textbf{VisReg JEPA (Fine-tuned)} & 0.8358 & 0.7215 & 12.19 px & 16.92 & \textbf{0.0132} & 18.21 ms/slice & 20.23 s/epoch \\
\bottomrule
\end{tabular}%
}
\end{table}

\subsection{Low-Data Label Efficiency Benchmark}
Table~\ref{tab:low_data_results} reports fine-tuning accuracy across 6 label availability tiers ($1\%$ to $100\%$ labels). At 10\% labels ($633$ slices), **SigReg JEPA** outperforms **nnU-Net** ($0.5198$ vs $0.5163$ Dice) with $2.4\times$ lower Hausdorff boundary distance ($16.44\,\text{px}$ vs $40.26\,\text{px}$). At 5\% labels ($316$ slices), SigReg JEPA ($0.2594$ Dice) outperforms standard I-JEPA ($0.0208$ Dice) by $>12\times$.

\begin{table}[ht]
\centering
\caption{Low-Data Label Efficiency Benchmark across 6 label availability fractions ($1\%$ to $100\%$).}
\label{tab:low_data_results}
\vspace{0.5em}
\resizebox{\textwidth}{!}{%
\begin{tabular}{lcccccc}
\toprule
\textbf{Model Architecture} & \textbf{1\% Labels (63 Slices)} & \textbf{5\% Labels (316 Slices)} & \textbf{10\% Labels (633 Slices)} & \textbf{25\% Labels (1.5k Slices)} & \textbf{50\% Labels (3.1k Slices)} & \textbf{100\% Labels (6.3k Slices)} \\
\midrule
\textbf{UNet Baseline} & 0.0286 & 0.0333 & 0.0486 & 0.0359 & 0.0377 & 0.0381 \\
\textbf{nnU-Net Baseline (SOTA)} & \textbf{0.1814} & \textbf{0.3945} & 0.5163 & \textbf{0.6045} & \textbf{0.6231} & 0.6256 \\
\textbf{I-JEPA (Fine-tuned)} & 0.0040 & 0.0208 & 0.3925 & 0.5934 & 0.5712 & 0.6410 \\
\textbf{SigReg JEPA (Fine-tuned)} & 0.0296 & \textbf{0.2594} & \textbf{0.5198} & 0.5893 & 0.6086 & \textbf{0.6471} \\
\textbf{VisReg JEPA (Fine-tuned)} & 0.0515 & 0.0341 & 0.3932 & 0.4766 & 0.5661 & 0.6122 \\
\bottomrule
\end{tabular}%
}
\end{table}

\subsection{Cross-Pathology \& Missing-Modality OOD Benchmark (\texttt{BraTS-MEN-RT})}
Table~\ref{tab:men_rt_results} evaluates zero-shot transfer on Extra-axial Meningioma MRI (\texttt{BraTS-MEN-RT}) providing only a single T1c sequence. Under zero-padded missing channels ($[0, \text{T1c}, 0, 0]$), standard **UNet completely collapses ($0.0078$ Dice)**, whereas **SigReg JEPA** achieves top zero-shot transfer ($0.1026$ Dice), outperforming nnU-Net ($0.0864$ Dice).

\begin{table}[ht]
\centering
\caption{Zero-Shot Cross-Pathology \& Missing-Modality OOD evaluation on BraTS-MEN-RT Meningioma MRI.}
\label{tab:men_rt_results}
\vspace{0.5em}
\resizebox{\textwidth}{!}{%
\begin{tabular}{llccc}
\toprule
\textbf{Adaptation Strategy} & \textbf{Model Architecture} & \textbf{Meningioma Test Dice} $\uparrow$ & \textbf{Test IoU} $\uparrow$ & \textbf{HD95 (px)} $\downarrow$ \\
\midrule
\textbf{Zero-Padding $[0, \text{T1c}, 0, 0]$} & \textbf{SigReg JEPA (Fine-tuned)} & \textbf{0.1026} & \textbf{0.0570} & \textbf{80.50 px} \\
Zero-Padding $[0, \text{T1c}, 0, 0]$ & nnU-Net Baseline (SOTA) & 0.0864 & 0.0461 & 111.49 px \\
Zero-Padding $[0, \text{T1c}, 0, 0]$ & VisReg JEPA (Fine-tuned) & 0.0596 & 0.0312 & 123.92 px \\
Zero-Padding $[0, \text{T1c}, 0, 0]$ & I-JEPA (Fine-tuned) & 0.0539 & 0.0278 & 134.97 px \\
Zero-Padding $[0, \text{T1c}, 0, 0]$ & UNet Baseline & 0.0078 (Collapse) & 0.0040 & 134.44 px \\
\midrule
\textbf{Channel Replication $[\text{T1c}, \text{T1c}, \text{T1c}, \text{T1c}]$} & nnU-Net Baseline (SOTA) & \textbf{0.1383} & \textbf{0.0750} & \textbf{75.20 px} \\
Channel Replication $[\text{T1c}, \text{T1c}, \text{T1c}, \text{T1c}]$ & \textbf{SigReg JEPA (Fine-tuned)} & 0.1196 & 0.0641 & 87.84 px \\
Channel Replication $[\text{T1c}, \text{T1c}, \text{T1c}, \text{T1c}]$ & UNet Baseline & 0.1095 & 0.0582 & 77.23 px \\
\bottomrule
\end{tabular}%
}
\end{table}


\section{Discussion \& Conclusion}

Our empirical investigation establishes three core scientific findings:
\begin{enumerate}
    \item \textbf{Label Efficiency in Low-Data Regimes}: At 10\% labels ($633$ training slices), SigReg JEPA outperforms fully supervised nnU-Net ($0.5198$ vs $0.5163$ Dice) while achieving $2.4\times$ sharper tumor boundary predictions ($16.44\,\text{px}$ HD95).
    \item \textbf{Robustness to Missing Modalities}: Standard CNNs collapse when modalities are missing at test time ($0.0078$ Dice for UNet). In contrast, SigReg JEPA pre-trained with variance-covariance regularization and Modality Dropout maintains top zero-shot cross-pathology transfer ($0.1026$ Dice).
    \item \textbf{Heuristic-Free Non-Collapse}: SigReg JEPA achieves full feature rank ($\mathbf{366.91}$), eliminating momentum teacher update overhead while running $\sim 40\%$ faster per epoch than nnU-Net.
\end{enumerate}

In conclusion, self-supervised JEPA models regularized via variance-covariance penalties offer a fast, label-efficient, and highly robust paradigm for multi-modal medical image segmentation.


\begin{ack}
This research was supported by the Master's Thesis Laboratory and computational hardware resources.
\end{ack}


\bibliographystyle{plainnat}
\bibliography{references}

\end{document}
